\documentclass[11pt]{article}
\usepackage[margin=1in]{geometry}
\usepackage{amsmath,amssymb,amsthm,mathtools}
\usepackage{graphicx}
\usepackage{xcolor}
\usepackage{enumitem}
\usepackage{microtype}
\usepackage{hyperref}
\hypersetup{colorlinks=true,linkcolor=blue!55!black,citecolor=blue!55!black,urlcolor=blue!55!black}

\newtheorem{theorem}{Theorem}
\newtheorem{lemma}{Lemma}
\newtheorem{corollary}{Corollary}
\newcommand{\Qp}{Q_{+}}
\newcommand{\Op}{\Omega_{+}}
\newcommand{\conv}{\operatorname{conv}}
\newcommand{\extr}{\operatorname{extr}}

\title{\vspace{-2em}Transitivity of Positive Extreme-Point Hulls\\
for Aligned Rearrangements}
\author{Solution packet for arXiv:1001.2426}
\date{17 August 2026}

\begin{document}
\maketitle

\begin{center}
\fcolorbox{blue!55!black}{blue!4}{\parbox{0.90\textwidth}{
\textbf{Status: likely valid substantial partial result.}
The general question in Remark~18 remains open.  The new theorem proves
transitivity whenever the intermediate rearrangement is obtained from the
initial rearrangement by a nonincreasing multiplier in $[0,1]$.  A separate
corollary records the full zero-dilation subcases already implicit in the
source paper.
}}
\end{center}

\section{Source and exact question}

The source is F.~Sukochev and D.~Zanin, \emph{Orbits in symmetric spaces},
J.\ Funct.\ Anal.\ \textbf{257} (2009), no.~1, 194--218,
arXiv:1001.2426 \cite{SZ}.  On page~12 of the compiled source PDF,
immediately after Lemma~17, the authors ask:

\begin{quote}
\emph{The collection of sets $\{Q(x),\,x\in E\}$ also satisfies the
transitivity property expressed in Lemma~17.  We do not know whether this is
the case for the collection $\{\Qp(x),\,x\in E\}$.}
\end{quote}

\begin{figure}[ht]
\centering
\includegraphics[width=0.96\textwidth]{figures/open_problem_crop.png}
\caption{Remark~18 in the source paper, rendered from the local source PDF.}
\end{figure}

Thus the question is whether
\[
        y\in \Qp(z),\quad z\in\Qp(x)
        \qquad\Longrightarrow\qquad y\in\Qp(x)
\]
for nonnegative elements of a symmetric Banach function space.

\section{Definitions}

Let $I=(0,1)$ or $(0,\infty)$, let $E=E(I)$ be a fully symmetric Banach
function space, and let $f^*$ denote the nonincreasing rearrangement.  For
$x\geq0$,
\[
 \Op(x)=\{u\geq0:u\prec\!\prec x\},\qquad
 \Qp(x)=\overline{\conv}^{\,\|\cdot\|_E}\extr\Op(x).
\]
The source proves that on $(0,1)$ the extreme rearrangements are
$x^*\mathbf 1_{[0,\beta]}$.  On $(0,\infty)$ these remain all the extreme
rearrangements when $x^*(\infty)=0$; if $x^*(\infty)>0$, there is also the
full-rearrangement possibility described in Corollary~31 of the source.

We use the standard rearrangement-transfer fact on a nonatomic Lebesgue
interval: a measure-preserving pullback preserves distributions, the norm of
a symmetric space, and the extreme-point descriptions above.  Consequently
$\Qp(x)$ depends only on $x^*$ and is invariant under equimeasurability.  This
is the same transfer principle used in the proof of Lemma~17 of the source.

\section{Proof intuition}

The useful feature of a nonincreasing multiplier $a\in[0,1]$ is that its
dyadic lower approximants are averages of indicators of initial intervals.
Multiplying by $x^*$ therefore turns each approximant into an average of
extreme points of $\Op(x)$.  The approximation is uniform in the multiplier,
so it converges in \emph{every} symmetric norm; no order continuity is needed.

If $z^*=a x^*$, every extreme point of $\Op(z)$ is obtained by one more
initial-interval cut (or by the allowed full rearrangement).  Its multiplier
relative to $x^*$ is still nonincreasing and lies in $[0,1]$.  Hence all
generators of $\Qp(z)$ already lie in $\Qp(x)$, which proves the desired
inclusion for this aligned class.

\section{The aligned-multiplier theorem}

\begin{lemma}[Monotone multipliers lie in the positive hull]\label{lem:mult}
Let $x\geq0$ belong to $E$, and let $b:I\to[0,1]$ be measurable and
nonincreasing.  If $w\geq0$ satisfies
\[
                  w^*=b\,x^*,
\]
then $w\in\Qp(x)$.
\end{lemma}

\begin{proof}
By rearrangement invariance it suffices first to treat $w=w^*=b x^*$.
For $n\geq1$ put
\[
 b_n=2^{-n}\lfloor 2^n b\rfloor
     =2^{-n}\sum_{j=1}^{2^n}
       \mathbf 1_{\{b\geq j2^{-n}\}}.
\]
Every level set of the nonincreasing function $b$ is, up to a null set, an
initial interval $[0,\beta_{n,j}]$.  Therefore
\[
 b_nx^*
 =2^{-n}\sum_{j=1}^{2^n}
   x^*\mathbf 1_{[0,\beta_{n,j}]}
\]
is a convex combination of extreme points of $\Op(x)$, using the
extreme-point characterization in the Appendix of \cite{SZ}.  The endpoints
$\beta=0$ and $\beta=\sup I$ cause no problem; in the semi-axis case $x^*$
itself satisfies the source's full-rearrangement tail condition.

Since $0\leq b-b_n<2^{-n}$, the ideal property of $E$ gives
\[
 \|(b-b_n)x^*\|_E\leq 2^{-n}\|x\|_E\longrightarrow0.
\]
Thus $b x^*\in\Qp(x)$.  Finally, applying the standard
measure-preserving rearrangement transfer to the approximating convex
combinations gives the same conclusion for every $w$ with
$w^*=b x^*$.
\end{proof}

\begin{theorem}[Transitivity for aligned rearrangements]\label{thm:main}
Let $E=E(I)$ be fully symmetric, let $x,z\in E$ be nonnegative, and suppose
that
\[
                 z^*(t)=a(t)x^*(t)\quad\text{a.e.}
\]
for some nonincreasing measurable $a:I\to[0,1]$.  Then
\[
                 z\in\Qp(x)
                 \qquad\text{and}\qquad
                 \Qp(z)\subseteq\Qp(x).
\]
In particular, $y\in\Qp(z)$ implies $y\in\Qp(x)$.
\end{theorem}

\begin{proof}
The first assertion is Lemma~\ref{lem:mult} with $b=a$.

Let $v\in\extr\Op(z)$.  On $(0,1)$, and on $(0,\infty)$ for every truncated
extreme point, the source's characterization gives
\[
 v^*=z^*\mathbf 1_{[0,\gamma]}
     =\bigl(a\mathbf 1_{[0,\gamma]}\bigr)x^*
\]
for some $\gamma$.  The multiplier
$a\mathbf 1_{[0,\gamma]}$ is again nonincreasing and takes values in
$[0,1]$, so Lemma~\ref{lem:mult} yields $v\in\Qp(x)$.

In the additional full-rearrangement case on $(0,\infty)$ one has
$v^*=z^*=a x^*$, and the same lemma again gives $v\in\Qp(x)$.
Thus every extreme point of $\Op(z)$ belongs to the closed convex set
$\Qp(x)$.  Taking closed convex hulls proves
$\Qp(z)\subseteq\Qp(x)$.
\end{proof}

\section{A full subcase already implicit in the source}

\begin{corollary}[Zero-dilation generators]\label{cor:zero}
If $\Qp(x)=\Op(x)$, then $\Qp(z)\subseteq\Qp(x)$ for every
$z\in\Qp(x)$.
In particular, the transitivity question has an affirmative answer whenever
one of the following source hypotheses holds:
\begin{enumerate}[label=(\roman*)]
\item $I=(0,1)$ and $\varphi(x)=0$;
\item $I=(0,\infty)$, $E\subseteq L_1$, and $\varphi_{\mathrm{fin}}(x)=0$;
\item $I=(0,\infty)$, $E\not\subseteq L_1$, and $\varphi(x)=0$.
\end{enumerate}
\end{corollary}

\begin{proof}
If $z\in\Qp(x)=\Op(x)$ and $y\in\Qp(z)\subseteq\Op(z)$, then
$y\prec\!\prec z\prec\!\prec x$.  Transitivity of submajorization gives
$y\in\Op(x)=\Qp(x)$.  The three displayed hypotheses imply
$\Qp(x)=\Op(x)$ by Theorems~23, 25, and 27 of \cite{SZ}, respectively.
\end{proof}

\section{Verification, novelty, and limitations}

\paragraph{Internal checks.}
The dyadic identity has exact coefficients summing to one.  The only norm
estimate is the lattice domination
$|(b-b_n)x^*|\leq2^{-n}x^*$, so the proof does not silently assume a Fatou
norm, separability, or order continuity.  The semi-axis proof separately
checks the extra full-rearrangement extreme points.

\paragraph{Bounded novelty check.}
The run registry, solution, attempt, and proof-gap indexes were searched for
arXiv:1001.2426 and the $Q_+$ transitivity language.  Exact-phrase and
core-keyword searches of arXiv/web results, together with inspection of the
closely related follow-up arXiv:1003.1817, found no explicit later resolution
of Remark~18 and no statement of Theorem~\ref{thm:main}.  This is a bounded
search, so novelty confidence is moderate rather than definitive.

\paragraph{Limitation and exact obstruction.}
The theorem does not cover an arbitrary $z\in\Qp(x)$.  Even
$z\prec\!\prec x$ need not imply $z^*\leq x^*$: mass may move from an early
high plateau into a later region where $x^*$ vanishes.  In that situation
$z^*/x^*$ is not an admissible decreasing multiplier.  Attempts to substitute
convex approximations are blocked because initial-interval truncation is
nonlinear and rearrangement of a sum need not preserve the aligned
decomposition.  This is precisely the remaining gap in the full problem.

\paragraph{Human-review recommendation.}
Check the rearrangement-transfer step in Lemma~\ref{lem:mult} against the
preferred convention for measure-preserving transformations on the
semi-axis, and verify that the full-tail extreme case is read exactly as in
Corollary~31 of the source.  The dyadic convex-hull mechanism itself is
elementary and appears robust.

\begin{thebibliography}{9}
\bibitem{SZ}
F.~Sukochev and D.~Zanin,
\emph{Orbits in symmetric spaces},
J.\ Funct.\ Anal.\ \textbf{257} (2009), no.~1, 194--218;
\href{https://arxiv.org/abs/1001.2426}{arXiv:1001.2426}.

\bibitem{KPS}
S.~G. Krein, Yu.~I. Petunin, and E.~M. Semenov,
\emph{Interpolation of Linear Operators},
Translations of Mathematical Monographs 54, American Mathematical Society,
1982.
\end{thebibliography}

\end{document}
